\chapter{Architectural Approach}

\section{Chosen Framework}

The selected architectural design framework is \textbf{Attribute-Driven Design (ADD)}.

ADD fits this project because the VerdeMart scenario is mainly driven by measurable quality
attribute pressures. The main question is not only which components should exist, but how the
commerce core should behave when surrounding operational systems are delayed, stale,
unavailable, or recovering.

In this project, ADD is applied as follows:

\begin{enumerate}
    \item Identify the business and architectural drivers for the omnichannel commerce core.
    \item Define quality attribute scenarios for the most important runtime pressures.
    \item Select architectural tactics that satisfy those scenarios.
    \item Decompose the target architecture into responsibilities, integration boundaries, and external adapters.
    \item Validate the design through a focused runtime demonstration of degradation and recovery.
\end{enumerate}

\section{Justification}

The chosen scenario requires nopCommerce to evolve from a web storefront into the commerce
core of a wider operational ecosystem. This creates architectural pressure in five main areas:
availability, consistency, performance, recoverability, and modifiability.

ADD is appropriate because each of these pressures can be expressed as a concrete quality
attribute scenario. Warehouse failure drives the need for asynchronous fulfillment and retry.
Stale inventory drives the need for inventory visibility, revalidation, and reconciliation. Order
bursts drive the need to protect customer-facing paths from slow backend synchronization.
Shipping outage drives the need for durable retry and explicit recovery state. New operational
channels drive the need for adapters and stable integration boundaries.

This gives a direct trace from the problem to the architecture:

\begin{center}
\begin{tabular}{|p{0.28\textwidth}|p{0.58\textwidth}|}
\hline
\textbf{ADD input} & \textbf{Architectural consequence} \\
\hline
Availability scenario & Avoid synchronous dependency on warehouse systems during checkout. \\
\hline
Consistency scenario & Introduce explicit inventory state, stale-state visibility, and reconciliation. \\
\hline
Performance scenario & Keep checkout and order reads separate from slow integration processing. \\
\hline
Recoverability scenario & Store failed integration work durably and retry after dependency recovery. \\
\hline
Modifiability scenario & Add external systems through adapters instead of direct database coupling. \\
\hline
\end{tabular}
\end{center}

\section{Why Not ACDM or ADM as the Main Framework}

\textbf{ACDM} was considered because it is useful for architectural decision-making, risk
discussion, and stakeholder validation. However, it is not the main framework for this project
because the strongest design input is the set of measurable quality attribute scenarios. The
project must show runtime behavior under failure, delay, stale state, and recovery; ADD gives a
more direct path from those scenarios to concrete architectural tactics and component
responsibilities.

ACDM-style reasoning is still useful later in the project, especially when documenting ADRs and
explaining rejected alternatives. For example, the decision to use asynchronous integration can
be recorded as an ADR with alternatives such as direct synchronous calls or shared database
access. However, the decision itself is primarily driven by the availability and recoverability
scenarios, which makes ADD the better primary method.

\textbf{ADM/TOGAF} was also considered, but it is broader than necessary for this assignment.
ADM is better suited to enterprise-wide transformation across business, data, application, and
technology architecture. VerdeMart does have an omnichannel business context, but the scope of
this work is a focused software architecture evolution of nopCommerce, not a full enterprise
architecture programme.

\section{How ADD Shapes the Target Architecture}

Using ADD means that each architectural element in the target design should answer a specific
quality attribute pressure. The surrounding systems are not added only because they exist in the
business scenario; they are added because they create design pressure that nopCommerce must
handle explicitly.

For this reason, the target architecture will emphasize asynchronous integration between
nopCommerce and operational systems. It will also make status explicit for pending, degraded,
failed, and recovered work, so that customers, support agents, and operators can understand
what is happening during dependency failures. Warehouse, inventory, and shipping flows will
use retry and reconciliation mechanisms instead of assuming that every external system is
always available. External systems will be connected through adapters to avoid coupling them
directly to nopCommerce internals, and data ownership will remain explicit so that surrounding
systems do not share the nopCommerce database.

The next chapter uses these ADD-driven conclusions to define the target architecture.
